\chapter{Large Language Models in Finance}\label{ch:ml:large-language-models-in-finance}

A \omterm{def:ml:large-language-models-in-finance:lm}{language model} asked whether headlines from its own training period were good news for the stock is right 70\% of the
time; asked about headlines written after its training data end, it is right 56\% of the time. The first number measures
its memory, the second its judgement, and a backtest that mixes the two measures nothing. The model in question is this
chapter's: a \omterm{def:ml:sequence-models-on-order-books:attention}{transformer} of 107\,000 parameters trained in a minute on the synthetic headlines of chapter~13, small enough to
be retrained with and without the evaluation period in its data. The protocol it tests is the one a desk needs for
commercial models whose training data it cannot see. The chapter then turns to what \omterm{def:ml:large-language-models-in-finance:lm}{language models} are used for in
practice, reading filings and news at scale through retrieval, and running research tasks as agents, and to what a desk can
and cannot hand them. Nothing in it calls a network service or downloads a model; every number comes from local code with
tests.

\section{Language models and how they are trained}

\begin{definition}[Language model, large language model]\label{def:ml:large-language-models-in-finance:lm}
A \emph{language model}\index{language model} assigns probabilities to sequences of tokens, usually through the conditional
probability of each token given the ones before it, $\P(w_t\mid w_1,\dots,w_{t-1})$; it generates text by sampling one
token at a time. A \emph{large language model}\index{large language model} (LLM) is a \omterm{def:ml:sequence-models-on-order-books:attention}{transformer} language model (chapter~8)
with billions of parameters, pre-trained (chapter~10) on a large share of the public text of the internet and books, then
fine-tuned to follow instructions.
\end{definition}

\begin{definition}[Prompt, context window, training cut-off]\label{def:ml:large-language-models-in-finance:prompt}
A \emph{prompt}\index{prompt} is the text given to a \omterm{def:ml:large-language-models-in-finance:lm}{language model}, which it continues. The \emph{context
window}\index{context window} is the largest number of tokens the model can condition on at once, prompt and answer
together. The \emph{training cut-off}\index{training cut-off} is the date of the latest text in the model's training data:
the model may know anything published before it and knows nothing published after it, except what the prompt tells it.
\end{definition}

\Cref{lst:ml:llm:tinylm} is the architecture at toy scale: token and position embeddings, two layers of causal
self-attention (a token attends only to earlier tokens), and a projection back onto the 505 words of the vocabulary. It is
trained, like its large relatives, to predict each next token. Its training text is the chapter-13 corpus, cut to its first
10\,000 headlines and written as sentences such as \texttt{w190 apex mining beat estimates amid copper demand => up}: a week,
the company, the headline, and after the arrow what the stock did (up or down, the sign of the reaction return). Asking the
model about a headline is prompting it with everything up to the arrow and comparing the probabilities it gives to
\texttt{up} and \texttt{down}. Two copies are trained: a \emph{clean} model whose training data end on day 907 and a
\emph{contaminated} model whose data run to day 1\,210; both are scored on the 5\,972 headlines before day 907, the 2\,016
between the two cut-offs and the 2\,012 after both.

Large models are built the same way at a scale this chapter cannot reach: BloombergGPT (Wu and co-authors, 2023), for
instance, is a 50-billion-parameter model trained on a 363-billion-token financial dataset built from Bloomberg's sources
together with 345 billion tokens of general text. What carries over from the toy is the protocol, not the numbers.

\section{Filings, calls and news at scale: retrieval}

A \omterm{def:ml:large-language-models-in-finance:lm}{language model} knows only its training data and its \omterm{def:ml:large-language-models-in-finance:prompt}{prompt}. To answer questions about a document collection that changes
every day (filings, transcripts, news), it is given the relevant passages in its \omterm{def:ml:large-language-models-in-finance:prompt}{prompt}.

\begin{definition}[Retrieval-augmented generation, BM25]\label{def:ml:large-language-models-in-finance:rag}
\emph{Retrieval-augmented generation}\index{retrieval-augmented generation} (RAG) answers a question by first retrieving
passages relevant to it from a document collection and then generating the answer from a \omterm{def:ml:large-language-models-in-finance:prompt}{prompt} that contains them (Lewis
and co-authors, 2020). \emph{BM25}\index{BM25} is the standard lexical retrieval score: for each query term found in a
passage, the term's inverse document frequency times a saturating function of its count in the passage, normalised by the
passage's length (Robertson and Zaragoza, 2009).
\end{definition}

The chapter's collection holds 540 passages from 135 synthetic annual filings, one passage per company, fiscal year and
metric (revenue, net debt, headcount, capital expenditure), each written with one of three synonyms (``revenue'', ``sales''
or ``turnover'') and two context words. One question is asked about each passage, worded with a synonym drawn at random, so
that about a third use the passage's own word; the answer is extracted from the top passage, so that a right passage
gives the right number. \Cref{fig:ml:llm:retrieval} compares four searches. \omterm{def:ml:large-language-models-in-finance:rag}{BM25} (\cref{lst:ml:llm:bm25}) finds every
passage asked about in its own words and 5.0\% of the others: a question about ``sales'' goes to another company's passage
that says ``sales''. Dense retrieval, the cosine between the average \omterm{def:ml:text-from-bag-of-words-to-embeddings:embedding}{word embeddings} of question and passage, knows that
``sales'' means revenue (the embeddings, trained on 5\,000 general business sentences, put the three synonyms at cosine
above 0.99) but not which company is asked about: 1.1\%. Adding the two standardised scores helps little (18.2\%). Expanding
the query with its words' embedding neighbours before running \omterm{def:ml:large-language-models-in-finance:rag}{BM25} finds 89.1\% of the paraphrased questions and 86.3\% of
the others: 88.1\% overall.

\begin{omfigure}
\begin{tikzpicture}
\begin{axis}[width=10.5cm,height=5.8cm,ybar,bar width=12pt,ymin=0,ymax=110,ylabel={right passage (\%)},area legend,
  xtick={0,1,2,3},xticklabels={BM25,dense,hybrid,expanded BM25},tick label style={font=\small},label style={font=\small},
  enlarge x limits=0.18,grid=major,grid style={gray!25},legend style={font=\footnotesize,at={(0.5,-0.2)},anchor=north,
  draw=none,fill=none},legend columns=2,table/col sep=comma]
\addplot[fill=omDef,draw=none] table[x=i,y=same]{figdata/ml/14-large-language-models-in-finance/retrieval.csv};
\addplot[fill=omThm,draw=none] table[x=i,y=other]{figdata/ml/14-large-language-models-in-finance/retrieval.csv};
\legend{question uses the passage's word,question uses a synonym}
\end{axis}
\end{tikzpicture}

\omcaption{Share of the 540 questions whose top passage is the one asked about, by search method and wording. Data:
\texttt{ml\_llm.retrieval}.}\label{fig:ml:llm:retrieval}
\end{omfigure}

\begin{definition}[Hallucination]\label{def:ml:large-language-models-in-finance:hallucination}
A \emph{hallucination}\index{hallucination} is a fluent answer that is not supported by the model's input or by fact: a
figure, a quotation or a reference produced because it is probable text, not because it is true.
\end{definition}

A retrieval system that always answers hallucinates by construction on questions its collection cannot answer. On 300
questions about company-years with no filing, the expanded search returns a passage and a number every time. A grounding
check, answering only when the passage names the question's company and fiscal year, abstains on all 300; on the answerable
questions it turns away 11.9\%, exactly the 64 of 540 whose top passage was about another company or year, which would have
been wrong. A generative model on top of retrieval needs the same discipline: quote the passage, check the entities, and
abstain rather than guess.

\section{Evaluation without leakage from the training cut-off}

\begin{definition}[Memorisation, anonymisation test]\label{def:ml:large-language-models-in-finance:memo}
\emph{Memorisation}\index{memorisation} is a model's ability to reproduce specific training examples, including their
outcomes, rather than a rule that generalises; \omterm{def:ml:large-language-models-in-finance:lm}{language models} memorise text seen even once (Carlini and co-authors, 2021).
The \emph{anonymisation test}\index{anonymisation test} replaces the identifiers in an input (company names, tickers, dates)
by placeholders and measures how much the model's accuracy falls: a model that reads the text loses little, a model that
recalls the event loses its memory.
\end{definition}

Evaluating a commercial \omterm{def:ml:large-language-models-in-finance:lm}{language model} on historical headlines is look-ahead bias (Book~7, chapter~3) in a new form: the
model may have read the headline and what the stock did next. \Cref{fig:ml:llm:contamination} measures both tests on the
chapter's two models. On the headlines before day 907, which both saw with their outcomes, the clean model is right 69.7\% of
the time and the contaminated one 65.4\% with names and weeks, and 57.7\% and 57.2\% without: that gap is memory. Between the two cut-offs the contaminated model is right
64.8\% of the time and the clean model 56.2\%, and anonymisation takes the contaminated model to 57.0\% and leaves the clean
one where it was (56.0\%). After both cut-offs the contaminated model is right 59.2\% of the time, the clean one 56.3\%, both
below the 61.6\% that the sign of the planted effect achieves.

\begin{omfigure}
\begin{tikzpicture}
\begin{axis}[width=11.5cm,height=6cm,ybar,bar width=8pt,ymin=50,ymax=72,ylabel={right direction (\%)},area legend,
  xtick={0,1,2},xticklabels={before both cut-offs,between the cut-offs,after both cut-offs},tick label style={font=\small},
  label style={font=\small},xmin=-0.5,xmax=2.5,grid=major,grid style={gray!25},legend columns=2,
  legend style={font=\footnotesize,at={(0.5,-0.2)},anchor=north,draw=none,fill=none},table/col sep=comma]
\addplot[fill=omDef,draw=none,bar shift=-13.5pt] table[x=i,y=clean]{figdata/ml/14-large-language-models-in-finance/contamination.csv};
\addplot[fill=omDef!40,draw=none,bar shift=-4.5pt] table[x=i,y=cleananon]{figdata/ml/14-large-language-models-in-finance/contamination.csv};
\addplot[fill=omThm,draw=none,bar shift=4.5pt] table[x=i,y=contam]{figdata/ml/14-large-language-models-in-finance/contamination.csv};
\addplot[fill=omThm!40,draw=none,bar shift=13.5pt] table[x=i,y=contamanon]{figdata/ml/14-large-language-models-in-finance/contamination.csv};
\addplot[black!70,dashed,thick,sharp plot,forget plot] table[x=x,y=y]{figdata/ml/14-large-language-models-in-finance/ceiling.csv};
\legend{clean,clean anonymised,contaminated,contaminated anonymised}
\end{axis}
\end{tikzpicture}

\omcaption{Share of headlines whose direction each model predicts correctly, by period, with names and weeks and with
placeholders. The clean model's training data end at the first cut-off, the contaminated model's at the second. Dashed: the
sign of the planted effect. Data: \texttt{ml\_llm.contamination}.}\label{fig:ml:llm:contamination}
\end{omfigure}

The chapter's named result follows. The \emph{\omterm{def:ml:large-language-models-in-finance:memo}{memorisation} premium}, the accuracy between the cut-offs minus the accuracy
after both, is 5.6 points for the contaminated model and $-0.1$ for the clean one; the \emph{anonymisation loss} between the
cut-offs is 7.8 points for the contaminated model and 0.2 for the clean one. Both tests catch the contamination without
knowing the training data, which is their purpose: a desk evaluating a model whose cut-off it doubts runs them, and
trusts only numbers from after the stated cut-off. The literature found the same two effects in real models. Lopez-Lira and
Tang (2023) evaluated GPT-4 on headlines published after its knowledge cut-off for that reason; Glasserman and Lin (2023)
compared original headlines with headlines stripped of company identifiers and found, in-sample, that the anonymised ones
did better: general knowledge of the company distracted the model more than look-ahead helped it. The toy has no general
knowledge to be distracted by, and after both cut-offs anonymisation costs each model about a point (59.2\% to 58.1\%,
56.3\% to 55.3\%): the \omterm{def:ml:large-language-models-in-finance:memo}{anonymisation test} is a test of memory, and its loss must be read against that small cost.

\begin{method}[Evaluating a language model on market data]\label{met:ml:llm:method}
\begin{enumerate}
\item Establish the \omterm{def:ml:large-language-models-in-finance:prompt}{training cut-off} from the provider's documentation and treat everything before it as in-sample.
\item Score only on text published after the cut-off; if history before it must be used, report the \omterm{def:ml:large-language-models-in-finance:memo}{anonymisation test} and
the accuracy by period, and expect them to disagree.
\item Freeze the model version, the \omterm{def:ml:large-language-models-in-finance:prompt}{prompt} and the decoding settings in the research log (Book~7, chapter~1): a provider's
silent update is a new model.
\item Compare against a cheap baseline (a dictionary, a tf-idf model) on the same headlines and horizon.
\end{enumerate}
\end{method}

\section{Agents in research workflows}

\begin{definition}[Language-model agent, tool call]\label{def:ml:large-language-models-in-finance:agent}
A \emph{language-model agent}\index{language-model agent} is a loop in which a \omterm{def:ml:large-language-models-in-finance:lm}{language model} chooses actions, observes their
results and decides the next action until a task is done (Yao and co-authors, 2023). A \emph{tool call}\index{tool call} is
one such action: the model emits a structured request (a function name and arguments) that the surrounding program executes,
returning the result to the model.
\end{definition}

An agent that can search, compute and run backtests is a research assistant with the same failure modes as a junior analyst,
at machine speed: it will use data it should not have. \Cref{lst:ml:llm:tools} is the firm's harness: tools are registered by
name, a dated tool receives the as-of date of the task and any result published after it is refused, and every call,
allowed or refused, goes to the log. The chapter's agent is a fixed plan rather than a \omterm{def:ml:large-language-models-in-finance:lm}{language model}, so that its behaviour
can be tested: asked, as of a random day, for the latest annual figure of a company, it searches, keeps the passages about the
company and the metric, and answers from the latest fiscal year it finds. On 400 such questions, without the guard, 80.8\%
of its answers come from filings published after the as-of date and 19.2\% are right; with the guard, none do and all are
right, and each question leaves one logged call.

\section{What a desk can and cannot delegate}

\omterm{def:ml:large-language-models-in-finance:lm}{Language models} are good at the work between data and decision that is expensive for people and cheap to check: extracting
fields from filings, summarising a call against the previous one, drafting code and tests, tagging news with events and
entities, answering questions over a collection with the passage quoted. Each of those has a checkable output. What cannot
be delegated is what cannot be checked cheaply: the decision to trade, the judgement that a backtest is clean, the claim that
a number is right when no passage supports it. Three rules keep the boundary. Nothing the model returns enters a
position or a report without a source that a person or a test can verify. Material non-public information (Book~7,
chapter~12) stays out of \omterm{def:ml:large-language-models-in-finance:prompt}{prompts} sent to services the firm does not control, since a \omterm{def:ml:large-language-models-in-finance:prompt}{prompt} is a disclosure. And every model
call that feeds research is logged with its version, \omterm{def:ml:large-language-models-in-finance:prompt}{prompt} and output, so that the result can be reproduced or at least
explained. Kim, Muhn and Nikolaev (2024) report that GPT-4, given standardised and anonymous financial statements, predicted
the direction of earnings changes better than analysts; the anonymisation is what made the result credible.

\section{Tutorial: the model that knew}\label{tut:ml:large-language-models-in-finance:tut}

\begin{tutorial}
\textbf{Goal.} Train a clean and a contaminated \omterm{def:ml:large-language-models-in-finance:lm}{language model} on the chapter-13 headlines and measure the \omterm{def:ml:large-language-models-in-finance:memo}{memorisation}
premium and the anonymisation loss; compare four retrieval methods and a grounding check; run the agent with and without its
guard. \textbf{End state:} \cref{fig:ml:llm:contamination,fig:ml:llm:retrieval}.
\begin{tutsteps}
\item \textbf{A decoder-only \omterm{def:ml:large-language-models-in-finance:lm}{language model} at toy scale.}
\omcode{code/firm/llmeval/firm_llmeval.py}{51}{67}{A causal transformer language model.}\label{lst:ml:llm:tinylm}
\item \textbf{\omterm{def:ml:large-language-models-in-finance:rag}{BM25}.}
\omcode{code/firm/llmeval/firm_llmeval.py}{122}{143}{The BM25 score.}\label{lst:ml:llm:bm25}
\item \textbf{Point-in-time tools.}
\omcode{code/firm/llmeval/firm_llmeval.py}{168}{192}{A tool registry with a point-in-time guard and a log.}\label{lst:ml:llm:tools}
\item \textbf{Run} \texttt{ml\_llm.contamination()} (about a minute on one core), \texttt{verdict()}, \texttt{retrieval()},
\texttt{unanswerable()}, \texttt{agent\_run()} and \texttt{fig\_llm.py}.
\end{tutsteps}
\textbf{What to change next.} Train the contaminated model for fewer \omterm{def:ml:neural-networks-for-noisy-tabular-data:backprop}{epochs} and watch the premium shrink; swap company names
between headlines (an entity-swap test) instead of removing them.
\end{tutorial}

\section{Build: evaluating language models without leakage}\label{bld:ml:large-language-models-in-finance:llmeval}

\begin{build}
\bfield{Purpose} \omterm{def:ml:large-language-models-in-finance:lm}{Language models} and retrieval used where their output can be checked, and evaluated only where their
training data cannot have seen the answer.
\bfield{Interface} \texttt{Vocab}, \texttt{TinyLM(vocab\_size, d, heads, layers, max\_len)}, \texttt{train\_lm(model, seqs,
epochs, lr, batch, seed)}, \texttt{next\_logprobs(model, prefixes)}; \texttt{BM25(docs, k1, b)}, \texttt{dense\_scores},
\texttt{expand\_query}, \texttt{hybrid}; \texttt{ToolRegistry(as\_of)} with \texttt{register(name, fn, dated)},
\texttt{call(name, **kw)} and \texttt{log}.
\bfield{Rules} No network calls and no downloaded weights in tests; scores reported by period relative to the \omterm{def:ml:large-language-models-in-finance:prompt}{training
cut-off}; every \omterm{def:ml:large-language-models-in-finance:agent}{tool call} logged; dated tools refuse anything after the as-of date.
\bfield{Acceptance tests} \texttt{code/firm/llmeval/tests/}: the \omterm{def:ml:large-language-models-in-finance:lm}{language model} memorises a small corpus and is causal (a
token's prediction does not change with later tokens); training is deterministic; \omterm{def:ml:large-language-models-in-finance:rag}{BM25} matches its formula on a hand example;
query expansion adds only close neighbours; the registry refuses a future document, logs every call and rejects unknown
tools.
\bfield{Stretch} A local open-weights model behind the same interface, run offline; an entity-swap test; a \omterm{def:ml:large-language-models-in-finance:prompt}{prompt} and version
registry in the experiment tracker (chapter~25).
\end{build}

\begin{omsources}
\item A.~Vaswani and co-authors, ``Attention is all you need'', arXiv:1706.03762, 2017.
\item P.~Lewis and co-authors, ``Retrieval-augmented generation for knowledge-intensive NLP tasks'', arXiv:2005.11401, 2020.
\item S.~Robertson and H.~Zaragoza, ``The probabilistic relevance framework: BM25 and beyond'', \emph{Foundations and Trends
in Information Retrieval} 3(4), 2009.
\item N.~Carlini and co-authors, ``Extracting training data from large language models'', arXiv:2012.07805, 2021.
\item A.~Lopez-Lira and Y.~Tang, ``Can ChatGPT forecast stock price movements? Return predictability and large language
models'', arXiv:2304.07619, 2023.
\item P.~Glasserman and C.~Lin, ``Assessing look-ahead bias in stock return predictions generated by GPT sentiment
analysis'', arXiv:2309.17322, 2023.
\item S.~Wu and co-authors, ``BloombergGPT: a large language model for finance'', arXiv:2303.17564, 2023.
\item S.~Yao and co-authors, ``ReAct: synergizing reasoning and acting in language models'', arXiv:2210.03629, 2023.
\item A.~Kim, M.~Muhn and V.~Nikolaev, ``Financial statement analysis with large language models'', arXiv:2407.17866, 2024.
\end{omsources}

\section{Exercises}

\begin{exercise}[$\star$]\label{exo:ml:large-language-models-in-finance:1}
A model's documentation gives a \omterm{def:ml:large-language-models-in-finance:prompt}{training cut-off} of 30 June. A backtest runs from January of the previous year to December.
Which part of it can measure the model's judgement, and what does the rest measure?
\end{exercise}

\begin{exercise}[$\star$]\label{exo:ml:large-language-models-in-finance:2}
Compute the \omterm{def:ml:large-language-models-in-finance:rag}{BM25} contribution of a term with document frequency 45 in 540 passages, appearing once in a passage of average
length, with $k_1 = 1.5$ and $b = 0.75$.
\end{exercise}

\begin{exercise}[$\star$]\label{exo:ml:large-language-models-in-finance:3}
From \cref{fig:ml:llm:contamination}, compute each model's \omterm{def:ml:large-language-models-in-finance:memo}{memorisation} premium and anonymisation loss.
\end{exercise}

\begin{exercise}[$\star\star$]\label{exo:ml:large-language-models-in-finance:4}
Why does dense retrieval alone find almost none of the right passages, and why does adding its score to \omterm{def:ml:large-language-models-in-finance:rag}{BM25}'s help so
little?
\end{exercise}

\begin{exercise}[$\star\star$]\label{exo:ml:large-language-models-in-finance:5}
Why is the contaminated model better than the clean one after both cut-offs (59.2\% against 56.3\%)? Is that evidence of
judgement?
\end{exercise}

\begin{exercise}[$\star\star$]\label{exo:ml:large-language-models-in-finance:6}
\emph{Find the flaw.} ``We asked the model for the direction of 10\,000 headlines from 2015 to 2023 and it was right 70\% of the
time, far above the 56\% of our tf-idf model; we removed the company names as a check and it was still 65\%.''
\end{exercise}

\begin{exercise}[$\star\star\star$]\label{exo:ml:large-language-models-in-finance:7}
\emph{Coding.} Verify that \texttt{TinyLM} is causal: change the last token of a \omterm{def:ml:large-language-models-in-finance:prompt}{prompt} and check that the model's
predictions at all earlier positions are unchanged. What in the code guarantees it?
\end{exercise}

\begin{exercise}[$\star\star\star$]\label{exo:ml:large-language-models-in-finance:8}
Design an agent task (research question, tools, as-of date, checks) for ``which of our signals decayed most last quarter?'',
and list what the log must contain for a reviewer to reproduce the answer.
\end{exercise}

\section{Problem: The Model That Knew}

\begin{problem}[{Weekend problem --- memory or judgement}]\label{pb:ml:large-language-models-in-finance:1}
The chapter's two \omterm{def:ml:large-language-models-in-finance:lm}{language models}, its filings collection and its agent.

\medskip\noindent\textbf{Part I --- The models.}
\begin{enumerate}
\item What does each model see in training, and how is it asked about a headline?
\item How large is the model and its vocabulary, and why does the protocol matter more than the size?
\item What is the best accuracy available after both cut-offs, and why is it not higher?
\item Why is one training headline in ten anonymised?
\end{enumerate}

\medskip\noindent\textbf{Part II --- Contamination.}
\begin{enumerate}[resume]
\item What are the two models' accuracies before both cut-offs, and what does the anonymised accuracy there show?
\item What are they between the cut-offs, and what explains the difference?
\item What does anonymisation do to each model between the cut-offs?
\item What does the literature say about the same two effects in real models?
\end{enumerate}

\medskip\noindent\textbf{Part III --- Retrieval and agents.}
\begin{enumerate}[resume]
\item How do the four searches do with the question in the passage's words and in other words?
\item What does the grounding check cost and what does it buy?
\item What does the agent do without its guard?
\item What should the log of an agent's run contain?
\end{enumerate}

\medskip\noindent\textbf{Part IV --- The verdict.}
\begin{enumerate}[resume]
\item State the \emph{named result}: the \omterm{def:ml:large-language-models-in-finance:memo}{memorisation} premium and the anonymisation loss for the contaminated and the clean
model.
\item How would you evaluate a vendor's news-sentiment model built on an LLM?
\item Which tasks would you give a \omterm{def:ml:large-language-models-in-finance:lm}{language model} first on a research desk?
\item What must never go into a \omterm{def:ml:large-language-models-in-finance:prompt}{prompt} sent outside the firm?
\item Why can an agent's backtest be worse than a person's?
\item How do you make an LLM-based research result reproducible?
\item When would you fine-tune a model rather than \omterm{def:ml:large-language-models-in-finance:prompt}{prompt} it?
\item In one sentence: what does a \omterm{def:ml:large-language-models-in-finance:lm}{language model}'s accuracy on history measure?
\end{enumerate}
\end{problem}

\section{Interview questions}

\begin{interviewq}[$\star$ \iqroles{researcher, mle}]\label{iq:ml:large-language-models-in-finance:1}
What is a \omterm{def:ml:large-language-models-in-finance:prompt}{training cut-off}, and why does it matter for a backtest that uses an LLM?
\end{interviewq}

\begin{interviewq}[$\star\star$ \iqroles{mle}]\label{iq:ml:large-language-models-in-finance:2}
Explain \omterm{def:ml:large-language-models-in-finance:rag}{retrieval-augmented generation}. How would you evaluate the retrieval step separately from the generation?
\end{interviewq}

\begin{interviewq}[$\star\star$ \iqroles{researcher}]\label{iq:ml:large-language-models-in-finance:3}
How would you test whether an LLM's apparent forecasting skill is \omterm{def:ml:large-language-models-in-finance:memo}{memorisation}?
\end{interviewq}

\begin{interviewq}[$\star\star$ \iqroles{mle, developer}]\label{iq:ml:large-language-models-in-finance:4}
\omterm{def:ml:large-language-models-in-finance:rag}{BM25} or embeddings for searching filings? Defend a design.
\end{interviewq}

\begin{interviewq}[$\star\star$ \iqroles{mle}]\label{iq:ml:large-language-models-in-finance:5}
How do you stop a research agent from using data after the as-of date, and how would you prove it did not?
\end{interviewq}

\begin{interviewq}[$\star\star\star$ \iqroles{researcher, trader}]\label{iq:ml:large-language-models-in-finance:6}
A vendor sells an LLM sentiment score with a backtested Sharpe ratio of 3 over ten years. What do you ask for before a trial?
\end{interviewq}
